\subsection{PSEUDO-REFLECTIONS}

DEFINITION 1. An endomorphism \(s\) of the vector space \(V\) is said to be a pseudo-reflection if \(1 - s\) is of rank 1.

Let \(s\) be a pseudo-reflection in \(V\) , and let \(D\) be the image of \(1 - s\) . By definition, \(D\) is of dimension 1; thus, given \(a \neq 0\) in \(D\) , there exists a non-zero linear form \(a^*\) on \(V\) such that \(x - s(x) = \langle x, a^* \rangle\) . \(a\) for all \(x \in V\) .

Conversely, given \(a \neq 0\) in V and a linear form \(a^{*} \neq 0\) on V, the formula

\[
s _ {a, a ^ {*}} (x) = x - \langle x, a ^ {*} \rangle . a \quad (x \in V)\tag{1}
\]

defines a pseudo-reflection \(s_{a,a^*}\) ; the image of \(1 - s_{a,a^*}\) is generated by \(a\) and the kernel of \(1 - s_{a,a^*}\) is the hyperplane of \(V\) consisting of those \(x\) such that \(\langle x, a^* \rangle = 0\) . If \(V^*\) is the dual of \(V\) , it is immediate that the transpose \(s_{a^*,a}\) of \(s_{a,a^*}\) is the pseudo-reflection of \(V^*\) given by the formula

\[
s _ {a ^ {*}, a} (x ^ {*}) = x ^ {*} - \langle x ^ {*}, a \rangle . a ^ {*} \quad (x ^ {*} \in V ^ {*}).\tag{2}
\]

If \(a\) is a non-zero vector, a pseudo-reflection with vector \(a\) is any pseudo-reflection \(s\) such that \(a\) belongs to the image of \(1 - s\) . The hyperplane of a pseudo-reflection \(s\) is the kernel of \(1 - s\) , the set of vectors \(x\) such that \(s(x) = x\) .

PROPOSITION 1. Let G be a group and \(\rho\) an irreducible linear representation of G on a vector space V; assume that there exists an element g of G such that \(\rho(g)\) is a pseudo-reflection.

\begin{enumerate}
\def\labelenumi{(\roman{enumi})}
\item
  Every endomorphism of V commuting with \(\rho(\mathbf{G})\) is a homothety, and \(\rho\) is absolutely irreducible.
\item
  Assume that V is finite dimensional. Let B be a non-zero bilinear form on V invariant under \(\rho(G)\) . Then B is non-degenerate, either symmetric or skew-symmetric, and every bilinear form on V invariant under \(\rho(G)\) is proportional to B.
\end{enumerate}

Let \(u\) be an endomorphism of \(V\) commuting with \(\rho(G)\) . Let \(g\) be an element of \(G\) such that \(\rho(g)\) is a pseudo-reflection and let \(D\) be the image of \(1 - \rho(g)\) . Since \(D\) is of dimension 1 and \(u(D) \subset D\) , there exists \(\alpha\) in \(K\) such that \(u - \alpha.1\) vanishes on \(D\) ; the kernel \(N\) of \(u - \alpha.1\) is then a vector subspace of \(V\) invariant under \(\rho(G)\) and is non-zero as it contains \(D\) ; since \(\rho\) is irreducible, \(N = V\) and \(u = \alpha.1\) . The second part of (i) follows from the first by Algebra, Chap. VIII, § 13, no. 4, Cor. of Prop. 5.

Let N (resp. \(N'\) ) be the subspace of V consisting of those x such that \(B(x,y)=0\) (resp. \(B(y,x)=0\) ) for all y in V; since B is invariant under \(\rho(G)\) , the subspaces N and \(N'\) of V are stable under \(\rho(G)\) and distinct from V since \(B\neq0\) . Since \(\rho\) is irreducible, \(N=N'=0\) and B is non-degenerate.

Since \(V\) is finite dimensional, every bilinear form on \(V\) is given by the formula

\[
\mathrm{B} ^ {\prime} (x, y) = \mathrm{B} (u (x), y)\tag{3}
\]

for some endomorphism u of V. If \(B'\) is invariant under \(\rho(G)\) , the endomorphism u commutes with \(\rho(G)\) . Indeed, let x, y be in V and let g be in G; since B and \(B'\) are invariant under \(\rho(G)\) , we have

\[
\begin{array}{r l} & {\mathrm{B} (u (\rho (g) (x)), y) = \mathrm{B} ^ {\prime} (\rho (g) (x), y) = \mathrm{B} ^ {\prime} (x, \rho (g ^ {- 1}) (y))} \\ & {\qquad = \mathrm{B} (u (x), \rho (g ^ {- 1}) (y)) = \mathrm{B} (\rho (g) (u (x)), y),} \end{array}
\]

hence \(u(\rho(g)(x)) = \rho(g)(u(x))\) since \(B\) is non-degenerate. By (i), there exists \(\alpha\) in \(K\) with \(u = \alpha.1\) , so \(B' = \alpha.B\) .

In particular, we can apply this to the bilinear form \(\mathrm{B}^{\prime}(x,y)=\mathrm{B}(y,x)\) ; then \(\mathrm{B}(y,x)=\alpha.\mathrm{B}(x,y)=\alpha^{2}.\mathrm{B}(y,x)\) for all x,y in V, and since B is nonzero we have \(\alpha^{2}=1\) , hence \(\alpha=1\) or \(\alpha=-1\) . Thus, B is either symmetric or skew-symmetric.

